% \section{Introduction} % no title as it's above main text

\IEEEPARstart{T}{ransformer-based} language models (LMs) are increasingly capable \cite{radford2019language, brown2020language, hoffmann2022empirical, chowdhery2023palm, kwa2025measuring} and widely used, including in highly sensitive domains like healthcare \cite{maity2025large}, law \cite{dehghani2025large} and cybersecurity \cite{ferrag2025generative}.
However, transformers have a black-box nature, meaning that while inputs and outputs are visible to humans, the internal process that connects them is opaque \cite{lipton2018mythos}.
Therefore, gaining a deeper understanding of the inner mechanisms of these models is essential to ensuring they behave safely and fairly by mitigating errors and biases, especially in critical and sensitive domains.

Mechanistic interpretability is an emerging field addressing this need by aiming to reverse-engineer neural networks such as transformers \cite{olah2020zoom,ferrando2024primer}.
For example, \citet{arditi2024refusal} identified a single activation vector in several safety and instruction-tuned LMs that mediated refusal behaviour, such that setting this vector to zero caused the model to stop refusing harmful requests, such as writing a fake defamatory article on a government official.
As such, a mechanistic understanding of neural networks is critical to understanding the limits of existing model safeguards and developing more robust ones \cite{yu2025robust}.

However, one major research gap in this field is understanding how transformers learn mechanisms for \emph{relational composition}, by which multiple vectors are composed into a single fixed-length vector \cite{wattenberg2024relational}. This enables transformers to perform well on tasks that require composing information from different input positions.
For example, \citet{brinkmann2024mechanistic} train a transformer to solve a pathfinding problem on binary trees. They represent the tree by its edges using sets of three tokens, $[\texttt{NODE}, \texttt{EDGE}, \texttt{SEP}]$, where \texttt{SEP} is a special token to separate different edges and does not contain any useful information for the actual task.
They find that models for which the tree depth exceeds the number of transformer layers learn to compute several subpaths in parallel, rather than building a single path one step per layer. They store these subpaths at \texttt{SEP} which acts as working memory between layers. In this way, the model composes several fixed-length token vectors into a single fixed-length activation vector at \texttt{SEP}. However, the authors do not investigate the structure of these activations or their links to relational composition.

\begin{figure*}[t]
  \centering
  \includegraphics[width=\linewidth]{figures/attn-flow.png}
  \caption{
    These graphs show the movement of information between token positions in our transformers on specific example inputs. 
    The three-layer model (right) implements a sequential copying algorithm, avoiding the relational composition task.
    We focus on a two-layer model (left) which has reduced accuracy (92\%) but performs the desired composition at the \texttt{SEP} token.
    Dotted lines represent the residual stream and red arrows indicate attention moving information between positions, with line thickness corresponding to attention magnitude. 
    As value and output matrices are frozen to the identity, attention performs scaled copying in the residual stream.
    Edges are annotated with the average impact on test accuracy when ablated (single-edge counterfactual); edges with no impact are not displayed. The numbers next to each node show the logit lens \citep{nostalgebraist2020logitlens-af, phukan-etal-2025-beyond} (the model's predicted probability distribution at that position and layer) intuitively demonstrating how the prediction distribution changes across layers.
  }
  \label{fig:attn-flow}
\end{figure*}

% Moreover, \citet{feng2023language} found an identity vector mechanism used in composing relation information in prompts such as ``Alice lives in France. Bob lives in Thailand.''. 

Whilst there is a substantial body of theoretical research into relational composition \citep{smolensky1990tpr, plate1995hrr, kanerva2009hdc, csordas-etal-2024-recurrent}, we lack empirical studies into how it actually manifests in LMs. 
\citet{wattenberg2024relational} propose additive matrix binding as a candidate mechanism, in which vectors are written to `slots' in the activation space using linear maps (distinct directions for each role). \citet{feng2023language} find broad use of binding ID vectors across families of LMs, which bind entities such as ``green'' and ``car'' in-context to represent ``green car'' (also using distinct directions as abstract identifiers).

These mechanisms have different implications for sparse autoencoders (SAEs): a popular method to recover interpretable features by decomposing dense, uninterpretable neuron activations into sparse, human-interpretable latents \cite{karvonen2025saebench, templeton2024scaling,mcdougall2025gemmascope2}.

SAEs rely on the assumption that neural network activations can be described in terms of independently interpretable features 
\citep{elhage2022toy}: each SAE latent should represent a single 
feature that can be understood in isolation. However, several 
composition mechanisms may violate this independence assumption in different ways. 
Additive matrix binding \citep{wattenberg2024relational} may produce ``echo features'': multiple latents that are copies of the same feature bound into different subspaces. ID vectors \citep{feng2023language}, meanwhile, may produce abstract latents whose interpretation cannot easily be grounded in the model's input. 

To bridge this research gap, we propose two research questions:

\begin{itemize}
    \item[RQ1] How does relational composition manifest in transformers?
    \item[RQ2] How do relational composition mechanisms impact SAEs?
\end{itemize}

For RQ1, we present a study of a toy model exhibiting another relational composition mechanism, which we call \emph{relative magnitude-based composition}: the model encodes relational information through the relative magnitudes of scalar attention weights applied to a fixed set of basis vectors, rather than by projecting information into distinct directions in the activation space. The same basis vectors represent both orderings of a bigram, $(x,y)$ and $(y,x)$, with order recoverable only from the ratio of their attention coefficients.
Toy models have been useful in the past for understanding phenomena that have then generalised to LMs, because they provide a controlled environment for study. For example, results on modulo arithmetic in toy models \citep{nanda2023progress} informed work on LMs \citep{baeumel2025modular}, and toy models of superposition \citep{elhage2022toy} led to the application of SAEs to LMs \citep{bricken2023monosemanticity}.
This part of our work was previously published in a workshop setting, see \cite{farrell2025order}.
For RQ2, we train a suite of SAEs on the toy model, analyse their results and discuss the significance for SAEs generally.
Our paper is structured as follows:
\begin{itemize}
    \item In Section \ref{s:method_rq1}, we present a two-layer attention-only transformer to encode two input tokens, $d_1$ and $d_2$, into a single separator token, \texttt{SEP}, and then reconstruct $d_1$ and $d_2$ as output. This design is motivated by the pathfinding model in \citet{brinkmann2024mechanistic}. 
    \item In Section \ref{s:method_rq2} we present our SAE training, evaluation and analysis methodology. We compare across BatchTopK \cite{bussmann2024batchtopk}, JumpReLU \cite{rajamanoharan2024jumping} and Matryoshka SAEs \cite{bussmann2025learning} to demonstrate the generality of our results to various state-of-the-art SAE configurations.
    \item In Section \ref{s:res_rq1}, we analyse the toy model and find that the order of the output tokens is defined by the relative magnitudes of scalar attention weights, rather than through a direction-based mechanism like additive matrix binding. This challenges the common perspective that SAE features can be interpreted as binary flags for feature presence \citep{quirke2025binary, paulo2024automatically}: in our model, the same features are active for both bigram orderings, $(a, b)$ and $(b, a)$, but with activation magnitudes that vary.
    \item In Section \ref{s:res_rq2}, we analyse the latents learned by the SAEs and perform latent steering experiments as in \citet{arditi2024refusal}. Distinct from additive matrix binding and ID vectors, we identify another failure mode for the independence assumption of SAE latents: latents whose meaning is contingent on co-active latents. We reliably identify latents within our highest-performing SAEs that represent the ordering relation between $d_1$ and $d_2$, but this is only interpretable given knowledge of what $d_1$ and $d_2$ are, which is information encoded in separate latents. No single latent is therefore independently sufficient; interpretation requires reasoning about latents jointly. These latents exhibit graded activations that can be reliably steered to reverse the original model outputs when the SAE's reconstruction is used instead of the original activation. We additionally provide a pipeline that, given an SAE trained on this model, reliably changes the order of between 25-70\% of the model's original outputs.
    \item In Section \ref{s:eval}, we discuss the implications of these results on the mechanistic interpretability agenda, and SAEs in particular.
\end{itemize}

We provide appendices in the report PDF in the code submission.
Our paper is intended as a foundational and exploratory work, and in future work we will attempt to validate our findings across less-constrained toy models and on pretrained LLMs. 
